\documentclass[11pt,letterpaper]{article}
\usepackage[T1]{fontenc}
\usepackage{amsmath,amssymb,amsthm,newtxtext,newtxmath,microtype,booktabs,array}
\usepackage[margin=1in]{geometry}
\usepackage[colorlinks=true,urlcolor=blue,linkcolor=black]{hyperref}
\newtheorem{proposition}{Proposition}
\newcommand{\Hh}{\mathcal H}
\newcommand{\E}{\mathbb E}
\newcommand{\C}{\mathcal C}
\newcommand{\Tr}{\operatorname{Tr}}
\newcommand{\SR}{\operatorname{SR}}
\setlength{\parskip}{4pt}
\setlength{\parindent}{0pt}
\begin{document}
\begin{center}
{\Large\bfseries Can kernel lengthscale determine portfolio shrinkage?}\\[5pt]
{Standalone mathematical note for \emph{The Law of Portfolio Learnability}}\\
10 October 2026
\end{center}

\section*{1. What is invariant, and what is not}
Write $M_t=N_t^{-1}\sum_iR^e_{i,t+1}\delta_{Z_{i,t}}$ for the random signed
measure of realized managed returns. A policy has payoff $R^p_{t+1}(W)=M_t(W)$.
Its RKHS Riesz representer is $X_{t,\ell}$, with
$\Sigma_\ell=\E[X_{t,\ell}\otimes X_{t,\ell}]$ and
$\mu_{\Hh_\ell}=\E X_{t,\ell}$. These objects include the return law and
cross-stock covariance; they are not the ordinary characteristic kernel operator.
Assume the same economy at every lengthscale and $\E\|X_{t,\ell}\|^2<\infty$.

\begin{proposition}[Conditional spectral order and transfer]
On a bounded full-dimensional Sobolev extension domain, fixed-smoothness
Euclidean Mat\'ern kernels at positive finite lengthscales have identical
RKHS sets and equivalent norms. If
\[
 a\|W\|_{\Hh_1}^2\leq\|W\|_{\Hh_2}^2\leq d\|W\|_{\Hh_1}^2,
 \qquad 0<a\leq d<\infty,
\]
then their positive managed eigenvalues satisfy
\[
 d^{-1}\mu_j(\ell_1)\leq\mu_j(\ell_2)\leq a^{-1}\mu_j(\ell_1).
\]
Thus an existing two-sided $j^{-b}$ order is preserved, but neither that
order nor a unique asymptotic amplitude is supplied by the kernel name.
If, additionally, $\mu_j(\ell)\sim c(\ell)j^{-b}$ with common $b>1$, the
asymptotic minimizer of $A(\ell)\lambda+B(\ell)\C_\ell(\lambda)/T$ obeys
\[
 \lambda_T^{\rm bound}(\ell)
 =\left[\frac{\rho(\ell)\kappa_b}{bT}c(\ell)^{1/b}\right]^{b/(b+1)},
 \quad\rho=B/A,\quad\kappa_b=\frac{\pi}{b\sin(\pi/b)}.
\]
This is a bound-proxy statement, not identification of a risk or Sharpe oracle.
\end{proposition}

\textit{Proof of the operator comparison.} The payoff map
$J_\ell:W\mapsto M_t(W)$ takes the same function set into $L^2$ of the
economy. Its squared singular values are the eigenvalues of $J_\ell^*J_\ell=\Sigma_\ell$.
The common quadratic form $\E[M_t(W)^2]$ has Rayleigh quotient divided by
$\|W\|_{\Hh_\ell}^2$. Apply the norm inequalities to every finite-dimensional
subspace and the compact-operator min--max principle. This yields the displayed
comparison, including equal null-space dimension. The norm equivalence follows
from the spectral-density ratio on page 2 and restriction with minimum extension
norm. The transfer calculation is proved on page 3. \hfill$\square$

The ordinary Sobolev correspondence requires appropriate domain regularity;
more general supports require separate trace/extension results. The 10,000-feature
empirical kernels are finite-dimensional approximations with changing cosine spans,
so this infinite-kernel proposition is not literally a theorem for their sample tails.

\newpage
\section*{2. Ordinary versus managed spectra and the amplitude}
Let $s=\nu+D/2$. With the Fourier convention $e^{2\pi i\omega\cdot z}$,
the unit-variance Euclidean Mat\'ern spectral density is
\[
 S_\ell(\omega)=\frac{2^D\pi^{D/2}\Gamma(s)(2\nu)^\nu}{\Gamma(\nu)\ell^{2\nu}}
 \left(\frac{2\nu}{\ell^2}+4\pi^2\|\omega\|^2\right)^{-s}.
\]
This normalization and the Gaussian comparison are given in
\href{https://gaussianprocess.org/gpml/chapters/RW4.pdf}{Rasmussen--Williams,
Chapter 4, equations (4.9)--(4.15)}. For $r=\ell_2/\ell_1\geq1$,
$r^{-2\nu}\leq S_{\ell_2}/S_{\ell_1}\leq r^D$. Fourier RKHS norms therefore
satisfy $r^{-D}\|W\|_1^2\leq\|W\|_2^2\leq r^{2\nu}\|W\|_1^2$.
Both spaces have the Sobolev order $s$, with lengthscale-dependent norms.
Restrictions to an extension domain preserve the same comparison constants.

Here is an explicit sufficient model for the nominal exponent. Periodize the
Euclidean covariance on a unit $D$-torus, \emph{without an extra diagonal
renormalization}. Its Fourier eigenvalues are $S_\ell(n)$, $n\in\mathbb Z^D$.
The elementary lattice count $\#\{n:\|n\|\leq R\}\sim v_DR^D$ gives
\[
 k_j(\ell)\sim d_\ell j^{-b},\qquad b=\frac{2s}{D}=1+\frac{2\nu}{D},
 \quad d_\ell\propto\ell^{-2\nu}.
\]
If a torus kernel is instead divided by its lengthscale-dependent diagonal,
the amplitude also divides by that normalizer. This distinction is explicit in
\href{https://arxiv.org/abs/2006.10160}{Borovitskiy et al. (2020)}; it cannot be
ignored when comparing kernel conventions.

To pass from $k_j$ to managed eigenvalues, suppose the economic second-moment
form on $L^2(P)$ is $\E[M_t(f)M_t(h)]=\langle f,Qh\rangle$ for a bounded
nonnegative operator $Q$. The managed operator is unitarily equivalent to
$K_\ell^{1/2}QK_\ell^{1/2}$, \emph{not} $K_\ell^2$.
If $q_-I\preceq Q\preceq q_+I$, min--max gives
$q_-k_j\leq\mu_j\leq q_+k_j$: the exponent transfers, but not a unique constant.
If $Q=qI+L$ with $q>0$ and compact $L$, then $\mu_j\sim qk_j$.
Indeed, approximate $L$ by finite-rank operators; the remaining form is bounded
by $\epsilon I$. Finite-rank interlacing changes rank by a fixed amount, while
regular variation makes $k_{j\pm m}/k_j\to1$. Then let $\epsilon\downarrow0$.

Independent characteristic draws and independent zero-mean idiosyncratic
returns with variance $\sigma^2$ give $Q=(\sigma^2/N)I$. Finite-rank factor
terms can be added if the idiosyncratic floor remains positive. A finite-factor
economy without that floor can have finite managed rank regardless of Mat\'ern
smoothness. Smooth return fields can instead contribute an additional smoothing
order. Varying $N$, return covariance, characteristic density, support geometry
or factor loadings requires a separate operator argument.
For a positive-mean example, constant mean $m$ and independent constant-variance
stock noise give $Q=((\sigma^2+m^2)/N)I+((N-1)m^2/N)\,1\otimes1$:
the same exponent survives a rank-one investment signal.

For $\nu=3/2$ and $D=130$, the nominal value is $133/130$, not $1+3/6$.
The repository's E5 assumes a managed spectral order directly; it does not prove
full-dimensional density, coercivity or a Weyl constant for the JKP economy.
No fitted ``effective dimension'' replaces 130 here.

\newpage
\section*{3. The exact asymptotic calculation, with its conditions}
Assume $\mu_j\sim cj^{-b}$, $c>0$, $b>1$. Put $J=(c/\lambda)^{1/b}$.
Riemann-sum comparison, with finite leading terms negligible relative to $J$,
gives
\[
 \C_\ell(\lambda)=\sum_j\frac{\mu_j}{\mu_j+\lambda}
 \sim J\int_0^\infty\frac{dx}{1+x^b}=\kappa_b(c/\lambda)^{1/b}.
\]
The integrability at infinity requires $b>1$. The beta integral evaluates to
$\pi/[b\sin(\pi/b)]$. Two-sided polynomial domination justifies discarding
the far tail, while pointwise asymptotic equivalence handles the growing bulk.
The sensitivity has the corresponding asymptotic
\[
 \sum_j\frac{\mu_j}{(\mu_j+\lambda)^2}
 \sim\frac{\kappa_b}{b}c^{1/b}\lambda^{-1-1/b}.
\]
This follows from the same argument with integrand $x^b/(1+x^b)^2$;
integration by parts evaluates its integral to $\kappa_b/b$.
Consequently the derivative equation for the approximate envelope is
$1=(\rho/T)(\kappa_b/b)c^{1/b}\lambda^{-1-1/b}$, giving Proposition 1.
Strict convexity makes this solution unique for a positive nonzero spectrum.

For the same $T$ and fixed common $b$, the corrected cross-lengthscale formulas are
\begin{align*}
 \frac{\lambda_2^{\rm bound}}{\lambda_1^{\rm bound}}
 &=\left(\frac{c_2}{c_1}\right)^{1/(b+1)}
   \left(\frac{\rho_2}{\rho_1}\right)^{b/(b+1)},\\
 \C_\ell(\lambda_\ell^{\rm bound})
 &\sim\kappa_b^{b/(b+1)}b^{1/(b+1)}
       \left(\frac{Tc_\ell}{\rho_\ell}\right)^{1/(b+1)},\\
 \frac{\C_2^{\rm bound}}{\C_1^{\rm bound}}
 &\sim\left(\frac{c_2}{c_1}\right)^{1/(b+1)}
       \left(\frac{\rho_2}{\rho_1}\right)^{-1/(b+1)}.
\end{align*}
The often proposed amplitude-only transfer is the special hypothesis
$\rho_2=\rho_1$. Under the additional torus/return assumptions on page 2,
$c_\ell\propto\ell^{-2\nu}$; then
$\lambda_2/\lambda_1=(\ell_2/\ell_1)^{-2\nu/(b+1)}$.
It is not an unconditional law of optimal portfolio shrinkage.

To make these statements uniform over a lengthscale family, require common $b$,
$c_\ell$ and $\rho_\ell$ bounded away from zero and infinity, uniform relative
tail equivalence, and uniformly summable dominating tails. A finite five-element
set permits taking maxima of finitely many tail thresholds. A continuum needs
genuine uniform assumptions. A family with $b\downarrow1$ needs separate care:
$\kappa_b$ diverges, convergence can be extremely slow, and uniform tail control
does not follow from pointwise statements. The numerical implementation evaluates
$\log\kappa_b$ stably near one; it never clips an invalid fitted $b$.

Constrained fixed penalty bounds can truncate these formulas. Admissibility and
the empirical stability regime of the paper's fixed-penalty learning bound must
also hold. Plugging a data-dependent spectrum into an argmin requires a
uniform-in-penalty theorem to inherit its guarantees. If $b$ itself changes,
both $\kappa_b$ and the exponents change; the displayed fixed-$b$ ratio is invalid.

\newpage
\section*{4. Why the spectrum does not identify an oracle constant}
Here is an exact \emph{Sharpe}-oracle counterexample for the actual response-one
ridge estimator. Let iid managed vectors take values
$\pm\sqrt2\,e_1$ and $\pm2\sqrt2\,e_2$, choosing each axis with probability
$1/2$. On a selected axis, a sign with mean $h\in(0,1)$ has probabilities
$(1+h)/2,(1-h)/2$. Economy A biases only the first axis; economy B biases only
the second. Both have $\Sigma=\operatorname{diag}(1,4)$ and identical sample
second-moment laws. Their means are $(m_1,0)$ and $(0,m_2)$, where
$m_1=h/\sqrt2$, $m_2=\sqrt2h$.

Fit ridge on $T=2$ observations. When each axis occurs once (probability $1/2$),
the absolute coefficient ratio is
\[
 r(\lambda)=\left|\frac{\widehat W_2}{\widehat W_1}\right|
 =\frac{2(1+\lambda)}{4+\lambda},\qquad r'(\lambda)>0.
\]
Integrating the \emph{population} Sharpe of each trained policy over the sixteen
ordered training outcomes gives
\begin{align*}
 \E\SR_A(\widehat W_\lambda)
 &=\frac{hm_1}{4\sqrt{1-m_1^2}}+
   \frac{hm_1}{2\sqrt{1-m_1^2+4r(\lambda)^2}},\\
 \E\SR_B(\widehat W_\lambda)
 &=\frac{hm_2}{4\sqrt{4-m_2^2}}+
   \frac{hm_2}{2\sqrt{4-m_2^2+r(\lambda)^{-2}}}.
\end{align*}
The first term in each case comes from drawing the signal axis twice. The
noise-axis-only payoff has zero population mean; exact zero policies are assigned
Sharpe zero. Conditional mean signs on the mixed-axis event are $h$, proving the
second terms. Thus A's expected Sharpe strictly decreases with $\lambda$, while
B's strictly increases. On any common positive finite interval, their unique
optimal penalties are opposite boundaries despite identical managed eigenvalues.
This controlled finite-law proof is not a simulated empirical portfolio result.

An analogous exact risk example has iid $X_t\in\{-1,+1\}$ with mean $m$:
its second moment is always one, but response-one expected-loss ridge has
$\lambda_Q^{\rm oracle}=(1-m^2)/(Tm^2)$. The tests verify both examples using
the existing estimator. Signal orientation and score noise cannot be recovered
from eigenvalues alone.

The paper's target has $\Sigma_\ell W_\ell^*=\mu_{\Hh_\ell}$. Its score is
$\zeta_{t,\ell}=(1-R^p_{t+1}(W_\ell^*))X_{t,\ell}$. Both the source norm and
the covariance of this score enter the learning proof. Changing Mat\'ern
lengthscale preserves the attainable policy set and hence its infimum population
loss, but changes the minimum RKHS norm and the regularization geometry.
If $R_\ell$ is the minimum norm among population minimizers, the page-2 bounds imply
\[
 r^{-D}R_1^2\leq R_2^2\leq r^{2\nu}R_1^2,\qquad
 r^{-2\nu}\|M_t\|_{\Hh_1^*}^2\leq\|M_t\|_{\Hh_2^*}^2
 \leq r^D\|M_t\|_{\Hh_1^*}^2.
\]
Thus a valid managed-tail scale can be transferred with
$B_{X,2}^2\leq r^D B_{X,1}^2$. The same economy's mixing coefficients are
unchanged, and the conditional fourth-moment condition concerns the same payoffs.
Under the paper's pricing condition its score bound is
$\E[\zeta\otimes\zeta]\preceq\sqrt{8K_4(1+K_4)}\,\Sigma_\ell$.
These are conservative comparisons, not formulas for sharp $A(\ell),B(\ell)$.
For $D=130$ the norm/tail bounds can be too wide to support useful transfer.

Moreover, an inequality admits many inflated choices of $A$ and $B$.
Unless their statistical functionals or a specific sharp bound are defined,
their ratio is not a unique parameter of the data-generating law. A one-time
CV anchor calibrates a tuning convention, not two identified structural constants.
The empirical ``implied rho'' diagnostics also require CV and cannot establish
a validation-free oracle method.

\newpage
\section*{5. Amplitude identification and the Gaussian negative control}
Two-sided bounds do not imply $j^b\mu_j\to c$. For example
\[
 \mu_j=j^{-b}\exp\{\epsilon\sin(\log\log(j+e))\},\qquad 0<\epsilon<b,
\]
is decreasing, positive and summable for $b>1$, and lies between
$e^{-\epsilon}j^{-b}$ and $e^{\epsilon}j^{-b}$, yet the normalized sequence
has no limit. Even regular variation with a slowly varying multiplier does not
force a convergent amplitude. A limiting principal symbol/Weyl law or explicit
relative asymptotic equivalence is stronger than the paper's two-sided E5.

An empirical operator has rank at most $\min(T,P)$. Its finite leading spectrum
can be extended by multiple decreasing summable infinite tails with different
asymptotic amplitudes. It cannot uniquely identify $c$ without tail assumptions.
For consistent amplitude estimation with known $b$, one could require growing
rank windows tending to infinity, uniform relative eigenvalue estimation errors
tending to zero on those windows, and uniform relative tail equivalence.
Then averages of $\log\widehat\mu_j+b\log j$ converge to $\log c$.
Fixed ranks 11--60 do not meet this argument; their intercept is a finite-range
statistic. Unknown $b$ additionally needs a suitably growing log-rank span and
control of approximation error. High regression $R^2$ alone supplies none of this.

A defensible \emph{regularized} historical diagnostic is possible. Fit a policy
on one past block and evaluate its score $(1-R^p(W))X$ on a disjoint later past
block. Under stationarity, summable dependence and sufficient moments, sample
averages and a consistent HAC estimator target that fixed regularized policy's
score covariance. They do not estimate the unregularized population source norm
or a sharp oracle penalty. Our annual fitted norm and score-trace aggregates are
explicitly descriptive in-sample versions of these regularized objects, not
plug-in estimates of $A,B$. Nonstationary economies can have identical entire
observed past laws and different future signal/noise; no past-only procedure can
identify both future oracles without additional stability assumptions.

The Gaussian density is proportional to
$\ell^D\exp(-2\pi^2\ell^2\|\omega\|^2)$ (the cited GPML normalization).
On the periodized model, lattice counting gives
\[
 \log k_j(\ell)=-2\pi^2\ell^2(j/v_D)^{2/D}+O(1)+o(j^{2/D}).
\]
The lengthscale changes the leading decay speed, not only an amplitude.
No polynomial Gaussian $b$ is imposed. Gaussian RKHS norms at different
lengthscales are generally not equivalent; their spaces can be strictly nested.
Finite-rank managed spectra and finite random features can obscure all these
population comparisons. This is why Gaussian is a useful empirical negative
control. Its explicitly requested transfer extension fits a finite-range surrogate,
without asserting a polynomial Gaussian population law.

\href{https://arxiv.org/abs/2309.10918}{Rosa et al. (2023)} give the relevant
Sobolev and restriction framework; they do not identify the JKP support dimension
or establish the managed-return coercivity required here.

\newpage
\section*{6. What the leading bound components can actually say}
One can bound the bias and score contribution without claiming sharp constants.
For the population ridge policy
$W_{\lambda,\ell}=(\Sigma_\ell+\lambda I)^{-1}\mu_{\Hh_\ell}$,
$W_{\lambda,\ell}-W_\ell^*=-\lambda(\Sigma_\ell+\lambda I)^{-1}W_\ell^*$.
The scalar inequality $\lambda^2x/(x+\lambda)^2\leq\lambda/4$ yields
\[
 \|\Sigma_\ell^{1/2}(W_{\lambda,\ell}-W_\ell^*)\|^2
 \leq\frac{\lambda}{4}\|W_\ell^*\|_{\Hh_\ell}^2.
\]
This gives a conservative \emph{population-bias} coefficient $R_\ell^2/4$.
It is not automatically the full finite-sample envelope's sharp $A(\ell)$.

Let $V_\ell=\E[\zeta_{t,\ell}\otimes\zeta_{t,\ell}]$ for the population
pricing score, with $V_\ell\preceq C_{\rm sc}\Sigma_\ell$ as on page 4.
Under iid observations or the paper's martingale-difference score condition,
the covariance of $\overline\zeta_T$ is $V_\ell/T$. A linearized fluctuation
using the \emph{population} regularized inverse consequently satisfies
\begin{align*}
 \E\|\Sigma_\ell^{1/2}(\Sigma_\ell+\lambda I)^{-1}
                 \overline\zeta_T\|^2
 &=\frac1T\Tr\{\Sigma_\ell(\Sigma_\ell+\lambda I)^{-1}
                      V_\ell(\Sigma_\ell+\lambda I)^{-1}\}\\
 &\leq\frac{C_{\rm sc}}{T}\sum_j
       \frac{\mu_j(\ell)^2}{(\mu_j(\ell)+\lambda)^2}
 \leq\frac{C_{\rm sc}}{T}\C_\ell(\lambda).
\end{align*}
Thus the same fourth-moment bound can supply a common conservative leading
score coefficient. Its sharp value depends on the orientation and size of
$V_\ell$, not only the eigenvalues of $\Sigma_\ell$.

Even these leading terms would require information about the source norm to
make their ratio specific. If one deliberately chooses this bias/score proxy,
its ratio is $4C_{\rm sc}/R_\ell^2$, and the page-4 norm comparisons bound its
change between $r^{-2\nu}$ and $r^D$. These are computable conservative
\emph{comparison} bounds given the stated proxy, not identification of the
unknown constants in a generic high-probability inequality. The actual estimator
also estimates its operator. Stability events, inverse-operator interaction,
tail scales, dependence and failure probability enter a full bound; the above
known-operator calculation cannot erase them or prove uniform data-dependent
selection guarantees. The repository's fixed-representation assumptions do not
identify their sharp lengthscale variation.

The regularized fitted norm and centered score-trace tables therefore target
observable finite-history diagnostics only. Recovering an unregularized source
norm needs a justified regularization limit and sufficient inverse-problem
control; a noisy minimum-norm estimate is not the population policy norm.
No fitted score trace is silently substituted for a structural $B(\ell)$.

\newpage
\section*{7. A testable algorithm, not a new optimality theorem}
For each kernel, recover only its baseline-lengthscale 1978 chronological
anchor: train on 1963--1972 formation months, validate on 1973--1977, and refit
on all 180 months. Define
\[
 \rho_0=\frac{180}{\sum_j\widehat\mu_{j,0}/(\widehat\mu_{j,0}+\lambda_0)^2}.
\]
Gaussian has $\lambda_0=2.515\times10^{-7}$, $\rho_0=4.344\times10^{-6}$;
Mat\'ern has $\lambda_0=1.559\times10^{-7}$, $\rho_0=1.590\times10^{-6}$.
Each experiment freezes its own single ratio. Sharing that ratio across its
five lengthscales is an explicit economic hypothesis.

\begin{center}\small
\begin{tabular}{p{.23\linewidth}p{.68\linewidth}}
\toprule Method & Penalty and lengthscale convention\\\midrule
Fixed lambda & Same raw $\lambda_0$, fixed baseline lengthscale.\\
Annual CV & Original 120-point grid, baseline lengthscale.\\
One-anchor spectral & Full past spectrum; unique proxy root, baseline lengthscale.\\
Joint CV & Original joint chronological lengthscale and penalty selection.\\
Spectral transfer & Fixed historical finite-range slope and robust amplitude proxy; lengthscale chosen from past performance. Gaussian is exploratory, with no population polynomial claim.\\
Adaptive full spectrum & Proxy root with frozen $\rho_0$; lengthscale chosen from past candidate performance.\\\bottomrule
\end{tabular}
\end{center}
The primary transfer formula is
\[
 \widehat\lambda_{t,\ell}=\lambda_0(180/T_t)^{b/(b+1)}
 \exp\!\left\{\frac{\widehat{\log c}_{t,\ell}-\widehat{\log c}_{0,\ell_0}}{b+1}\right\}.
\]
Each kernel's surrogate $b$ is frozen from its initial baseline ranks 11--60;
Gaussian has proxy slope 1.6234. Gaussian transfer is a documented user-requested
exploratory amendment after the initial Mat\'ern-primary run, not an imposed
population decay law. The unverified nominal $133/130$ is Mat\'ern-only;
$1.5,2,3$ are finite-range sensitivities.
Rank intervals, prefix counts, numerical thresholds and an amplitude-identification
gate are frozen before new results. A failed gate rejects a structural interpretation;
the mechanical transfer remains a labelled falsification exercise. No exponent is
clipped to make $b>1$. The adjusted-ratio variant is separate and requires five
1978 CV calibrations; boundary values are identified interval endpoints only.

All candidates are refit from past observations each year. Adaptive lengthscale
selection first occurs in 1984, using the last 60 genuinely OOS candidate returns
from decisions 1979--1983 and minimum response-one loss. Anchor validation
returns are never recycled as fresh candidate evidence. Thirty-six-month selection
is secondary. Penalty bounds remain each representation's original initial grid.
The 1984--2024 common sample has 492 returns, February 1984--January 2025.
All net accounts in that comparison start from cash together, charging entry;
the old full-history fixed accounts are separately reconciled.

Costs use each policy's actual signed weights and the repository's exact
self-financing 25-bp two-sided traded notional and 30-bp/year short fee.
Target-turnover subtraction is separately labelled. Paired circular blocks of
12 months, with 6 and 24 as sensitivities, use 5,000 shared draws conditional on
the fitted paths. Hindsight grid minima/maxima are diagnostics only, not
population oracles or calibration inputs. The initial Mat\'ern identification gate
fails; empirical performance and the final GO/NO-GO decision are reported
separately. Neither a smooth complexity path nor numerical similarity proves a
population spectral law or financial optimality.
\end{document}
